% --- Archon physics formalization source begin ---
% archon:physics
% archon:covers PhyXMiniProblems/problem_phyx_mini_0152.lean
% archon:source-report reports/phyx_mini/problem_phyx_mini_0152.source.json

\chapter{Physics problem phyx\_mini\_0152}
\label{ch:PhyXMiniProblems_problem_phyx_mini_0152}

\paragraph{Problem source.}
\#\# Physical scenario

One of the beams of an interferometer (figure) passes through a small evacuated glass container \textbackslash{}(1.155\textbackslash{},\textbackslash{}mathrm\{cm\}\textbackslash{}) deep. When a gas is allowed to slowly fill the container, a total of 158 dark fringes are counted to move past a reference line. The light used has a wavelength of \textbackslash{}(632.8\textbackslash{},\textbackslash{}mathrm\{nm\}\textbackslash{}).Assumed that the interferometer is in vacuum.

\#\# Question

Calculate the index of refraction of the gas at its final density.

\#\# Answer choices

- A: A: 1.004008

- B: B: 1.004828

- C: C: 1.004328

- D: D: 1.005018

\#\# Figure caption (auxiliary; use the image as primary evidence)

The image depicts an optical setup involving mirrors and a glass container. Here's a detailed description:

- **Source**: The starting point on the left where light originates.
- **Light Path**: Pink arrows indicate the direction of the light.
- **Mirror \textbackslash{}( M\_S \textbackslash{})**: Positioned at a 45-degree angle to the light path, reflecting light upwards.
- **Mirror \textbackslash{}( M\_1 \textbackslash{})**: Indicated to be above, receiving light reflected by \textbackslash{}( M\_S \textbackslash{}).
- **Glass Container**: Positioned on the horizontal path, with a label indicating its width as 1.155 cm.
- **Mirror \textbackslash{}( M\_2 \textbackslash{})**: Positioned on the right end of the light path, reflecting light back towards the source.

The setup likely demonstrates an interferometer or similar optical experiment using mirrors and a glass medium.

\#\# Dataset metadata

Wave Optics; Multi-Formula Reasoning, Physical Model Grounding Reasoning

\paragraph{Recorded answer/context.}
C: C: 1.004328

\paragraph{Figure/image path.}
/root/proposal\_for\_physic/science-mango/phyx\_mini\_run/phyx\_data/test\_image/152.png

\paragraph{Formalization target.}
create a compiling Lean file with sorry bodies at `PhyXMiniProblems/problem\_phyx\_mini\_0152.lean`. The Lean declarations must preserve the physical quantities, dimensions or dimensional roles, figure labels, governing-law hypotheses, and final relation expressed by this problem.
Use Mathlib/Physlib names found through LeanExplore where available. If a physics API is missing, introduce faithful local abstractions rather than scalar placeholder aliases.

\begin{theorem}[Physics formalization target]
\label{thm:physics:phyx_mini_0152:target}
\lean{PhyXMiniProblems.ProblemPhyXMini0152.problem_phyx_mini_0152}
\uses{def:physics:phyx-mini-0152:phyxminiproblems-problemphyxmini0152-gascellinterferometer, def:physics:phyx-mini-0152:phyxminiproblems-problemphyxmini0152-hasstatedreadouts, def:physics:phyx-mini-0152:phyxminiproblems-problemphyxmini0152-matchessourcefigure, def:physics:phyx-mini-0152:phyxminiproblems-problemphyxmini0152-hasphysicalparameters, def:physics:phyx-mini-0152:phyxminiproblems-problemphyxmini0152-hasvacuumreferenceindex, def:physics:phyx-mini-0152:phyxminiproblems-problemphyxmini0152-obeysgascellopticalpathlaw, def:physics:phyx-mini-0152:phyxminiproblems-problemphyxmini0152-obeysfringecountinglaw, def:physics:phyx-mini-0152:phyxminiproblems-problemphyxmini0152-answerrefractiveindex, def:physics:phyx-mini-0152:phyxminiproblems-problemphyxmini0152-isnearestmillionthreadout}
The assigned autoformalize agent should translate this physics problem into Lean declarations in the covered file, with theorem and lemma proof bodies written as `by sorry`.
\end{theorem}
\begin{proof}
This is an autoformalization task, not a proof task. Produce statements that can later be proved by the physics prover without weakening the physical model.
\end{proof}

% --- Archon declaration topology begin ---
\section{Declaration topology}
\begin{definition}[Length Quantity]
  \label{def:physics:phyx-mini-0152:phyxminiproblems-problemphyxmini0152-lengthquantity}
  \lean{PhyXMiniProblems.ProblemPhyXMini0152.LengthQuantity}
  A signed physical length with coherent readouts in every unit system
\end{definition}
\begin{proof}
  This is the declared model definition of Length Quantity.
\end{proof}
\begin{definition}[length In Meters]
  \label{def:physics:phyx-mini-0152:phyxminiproblems-problemphyxmini0152-lengthinmeters}
  \lean{PhyXMiniProblems.ProblemPhyXMini0152.lengthInMeters}
  \uses{def:physics:phyx-mini-0152:phyxminiproblems-problemphyxmini0152-lengthquantity}
  The SI scalar readout of a physical length, in metres
\end{definition}
\begin{proof}
  This is the declared model definition of length In Meters.
\end{proof}
\begin{definition}[length In Centimeters]
  \label{def:physics:phyx-mini-0152:phyxminiproblems-problemphyxmini0152-lengthincentimeters}
  \lean{PhyXMiniProblems.ProblemPhyXMini0152.lengthInCentimeters}
  \uses{def:physics:phyx-mini-0152:phyxminiproblems-problemphyxmini0152-lengthquantity}
  The scalar readout of a physical length, in centimetres
\end{definition}
\begin{proof}
  This is the declared model definition of length In Centimeters.
\end{proof}
\begin{definition}[length In Nanometers]
  \label{def:physics:phyx-mini-0152:phyxminiproblems-problemphyxmini0152-lengthinnanometers}
  \lean{PhyXMiniProblems.ProblemPhyXMini0152.lengthInNanometers}
  \uses{def:physics:phyx-mini-0152:phyxminiproblems-problemphyxmini0152-lengthquantity}
  The scalar readout of a physical length, in nanometres
\end{definition}
\begin{proof}
  This is the declared model definition of length In Nanometers.
\end{proof}
\begin{definition}[Cell State]
  \label{def:physics:phyx-mini-0152:phyxminiproblems-problemphyxmini0152-cellstate}
  \lean{PhyXMiniProblems.ProblemPhyXMini0152.CellState}
  The two cell states compared during the slowly performed experiment
\end{definition}
\begin{proof}
  This is the declared model definition of Cell State.
\end{proof}
\begin{definition}[Figure Element]
  \label{def:physics:phyx-mini-0152:phyxminiproblems-problemphyxmini0152-figureelement}
  \lean{PhyXMiniProblems.ProblemPhyXMini0152.FigureElement}
  Labels in the source figure, together with the stated observation line
\end{definition}
\begin{proof}
  This is the declared model definition of Figure Element.
\end{proof}
\begin{definition}[Beam Leg]
  \label{def:physics:phyx-mini-0152:phyxminiproblems-problemphyxmini0152-beamleg}
  \lean{PhyXMiniProblems.ProblemPhyXMini0152.BeamLeg}
  Directed pieces of the beam path visible in the source figure
\end{definition}
\begin{proof}
  This is the declared model definition of Beam Leg.
\end{proof}
\begin{definition}[Propagation Direction]
  \label{def:physics:phyx-mini-0152:phyxminiproblems-problemphyxmini0152-propagationdirection}
  \lean{PhyXMiniProblems.ProblemPhyXMini0152.PropagationDirection}
  Cardinal directions used to transcribe the arrows in the figure
\end{definition}
\begin{proof}
  This is the declared model definition of Propagation Direction.
\end{proof}
\begin{definition}[Gas Cell Interferometer]
  \label{def:physics:phyx-mini-0152:phyxminiproblems-problemphyxmini0152-gascellinterferometer}
  \lean{PhyXMiniProblems.ProblemPhyXMini0152.GasCellInterferometer}
  \uses{def:physics:phyx-mini-0152:phyxminiproblems-problemphyxmini0152-lengthquantity, def:physics:phyx-mini-0152:phyxminiproblems-problemphyxmini0152-cellstate, def:physics:phyx-mini-0152:phyxminiproblems-problemphyxmini0152-figureelement, def:physics:phyx-mini-0152:phyxminiproblems-problemphyxmini0152-beamleg, def:physics:phyx-mini-0152:phyxminiproblems-problemphyxmini0152-propagationdirection}
  This structure records the Gas Cell Interferometer component of the problem-specific physical model
\end{definition}
\begin{proof}
  This is the declared model definition of Gas Cell Interferometer.
\end{proof}
\begin{definition}[Has Stated Readouts]
  \label{def:physics:phyx-mini-0152:phyxminiproblems-problemphyxmini0152-hasstatedreadouts}
  \lean{PhyXMiniProblems.ProblemPhyXMini0152.HasStatedReadouts}
  \uses{def:physics:phyx-mini-0152:phyxminiproblems-problemphyxmini0152-lengthincentimeters, def:physics:phyx-mini-0152:phyxminiproblems-problemphyxmini0152-lengthinnanometers, def:physics:phyx-mini-0152:phyxminiproblems-problemphyxmini0152-gascellinterferometer}
  This def records the Has Stated Readouts component of the problem-specific physical model
\end{definition}
\begin{proof}
  This is the declared model definition of Has Stated Readouts.
\end{proof}
\begin{definition}[Matches Source Figure]
  \label{def:physics:phyx-mini-0152:phyxminiproblems-problemphyxmini0152-matchessourcefigure}
  \lean{PhyXMiniProblems.ProblemPhyXMini0152.MatchesSourceFigure}
  \uses{def:physics:phyx-mini-0152:phyxminiproblems-problemphyxmini0152-gascellinterferometer}
  This def records the Matches Source Figure component of the problem-specific physical model
\end{definition}
\begin{proof}
  This is the declared model definition of Matches Source Figure.
\end{proof}
\begin{definition}[Has Physical Parameters]
  \label{def:physics:phyx-mini-0152:phyxminiproblems-problemphyxmini0152-hasphysicalparameters}
  \lean{PhyXMiniProblems.ProblemPhyXMini0152.HasPhysicalParameters}
  \uses{def:physics:phyx-mini-0152:phyxminiproblems-problemphyxmini0152-lengthinmeters, def:physics:phyx-mini-0152:phyxminiproblems-problemphyxmini0152-gascellinterferometer}
  Positivity and monotonicity conditions for the measured physical system
\end{definition}
\begin{proof}
  This is the declared model definition of Has Physical Parameters.
\end{proof}
\begin{definition}[Has Vacuum Reference Index]
  \label{def:physics:phyx-mini-0152:phyxminiproblems-problemphyxmini0152-hasvacuumreferenceindex}
  \lean{PhyXMiniProblems.ProblemPhyXMini0152.HasVacuumReferenceIndex}
  \uses{def:physics:phyx-mini-0152:phyxminiproblems-problemphyxmini0152-gascellinterferometer}
  The initially evacuated cell has the vacuum refractive index 1
\end{definition}
\begin{proof}
  This is the declared model definition of Has Vacuum Reference Index.
\end{proof}
\begin{definition}[Obeys Gas Cell Optical Path Law]
  \label{def:physics:phyx-mini-0152:phyxminiproblems-problemphyxmini0152-obeysgascellopticalpathlaw}
  \lean{PhyXMiniProblems.ProblemPhyXMini0152.ObeysGasCellOpticalPathLaw}
  \uses{def:physics:phyx-mini-0152:phyxminiproblems-problemphyxmini0152-lengthinmeters, def:physics:phyx-mini-0152:phyxminiproblems-problemphyxmini0152-gascellinterferometer}
  This def records the Obeys Gas Cell Optical Path Law component of the problem-specific physical model
\end{definition}
\begin{proof}
  This is the declared model definition of Obeys Gas Cell Optical Path Law.
\end{proof}
\begin{definition}[Obeys Fringe Counting Law]
  \label{def:physics:phyx-mini-0152:phyxminiproblems-problemphyxmini0152-obeysfringecountinglaw}
  \lean{PhyXMiniProblems.ProblemPhyXMini0152.ObeysFringeCountingLaw}
  \uses{def:physics:phyx-mini-0152:phyxminiproblems-problemphyxmini0152-lengthinmeters, def:physics:phyx-mini-0152:phyxminiproblems-problemphyxmini0152-gascellinterferometer}
  This def records the Obeys Fringe Counting Law component of the problem-specific physical model
\end{definition}
\begin{proof}
  This is the declared model definition of Obeys Fringe Counting Law.
\end{proof}
\begin{definition}[Answer Choice]
  \label{def:physics:phyx-mini-0152:phyxminiproblems-problemphyxmini0152-answerchoice}
  \lean{PhyXMiniProblems.ProblemPhyXMini0152.AnswerChoice}
  Labels of the four numerical choices printed with the problem
\end{definition}
\begin{proof}
  This is the declared model definition of Answer Choice.
\end{proof}
\begin{definition}[answer Refractive Index]
  \label{def:physics:phyx-mini-0152:phyxminiproblems-problemphyxmini0152-answerrefractiveindex}
  \lean{PhyXMiniProblems.ProblemPhyXMini0152.answerRefractiveIndex}
  \uses{def:physics:phyx-mini-0152:phyxminiproblems-problemphyxmini0152-answerchoice}
  The displayed dimensionless refractive-index readout for each choice
\end{definition}
\begin{proof}
  This is the declared model definition of answer Refractive Index.
\end{proof}
\begin{definition}[Is Nearest Millionth Readout]
  \label{def:physics:phyx-mini-0152:phyxminiproblems-problemphyxmini0152-isnearestmillionthreadout}
  \lean{PhyXMiniProblems.ProblemPhyXMini0152.IsNearestMillionthReadout}
  This def records the Is Nearest Millionth Readout component of the problem-specific physical model
\end{definition}
\begin{proof}
  This is the declared model definition of Is Nearest Millionth Readout.
\end{proof}
\begin{lemma}[final Refractive Index eq exact]
  \label{lem:physics:phyx-mini-0152:phyxminiproblems-problemphyxmini0152-finalrefractiveindex-eq-exact}
  \lean{PhyXMiniProblems.ProblemPhyXMini0152.finalRefractiveIndex_eq_exact}
  \uses{def:physics:phyx-mini-0152:phyxminiproblems-problemphyxmini0152-gascellinterferometer, def:physics:phyx-mini-0152:phyxminiproblems-problemphyxmini0152-hasstatedreadouts, def:physics:phyx-mini-0152:phyxminiproblems-problemphyxmini0152-matchessourcefigure, def:physics:phyx-mini-0152:phyxminiproblems-problemphyxmini0152-hasphysicalparameters, def:physics:phyx-mini-0152:phyxminiproblems-problemphyxmini0152-hasvacuumreferenceindex, def:physics:phyx-mini-0152:phyxminiproblems-problemphyxmini0152-obeysgascellopticalpathlaw, def:physics:phyx-mini-0152:phyxminiproblems-problemphyxmini0152-obeysfringecountinglaw}
  This lemma records the final Refractive Index eq exact component of the problem-specific physical model
\end{lemma}
\begin{proof}
  The conclusion follows from the governing relations and figure readouts named in the statement of final Refractive Index eq exact.
\end{proof}
% --- Archon declaration topology end ---
% --- Archon physics formalization source end ---
